% answers_short.tex —— 2007 年·普通高等学校招生考试·陕西卷·文科 · 简短答案速查
% 由 tools/make_short_answers.py 自 answers.tex 的 \gksolhead/\gkanswer 提取，
% 并与 2007/2007_数学真题_陕西文_答案要点.md 逐题比对后生成，勿手改（改源 MD 后重跑）。

\begin{gknotice}
\textbf{使用说明}：本节只给最终答案便于快速核对，完整推导见解析版。答案为编写者独立推导，\textbf{非官方标准答案}。
\end{gknotice}

\gkpart{选择题}
\begin{center}\begin{tabular}{ccccccccccccc}
\toprule
题号 & 1 & 2 & 3 & 4 & 5 & 6 & 7 & 8 & 9 & 10 & 11 & 12 \\
\midrule
答案 & C & B & B & A & C & C & D & D & B & A & C & D \\
\bottomrule
\end{tabular}\end{center}

\gkpart{填空题}
\begin{itemize}[leftmargin=2.6em]
  \item[\textbf{13.}] $40$。
  \item[\textbf{14.}] $8$。
  \item[\textbf{15.}] $60$。
  \item[\textbf{16.}] $6$。
\end{itemize}

\gkpart{解答题}
\begin{itemize}[leftmargin=2.6em]
  \item[\textbf{17.}] （1）$m=1$；（2）最小值 $1-\sqrt2$。
  \item[\textbf{18.}] （1）$\dfrac{96}{625}$；（2）$\dfrac{101}{125}$。
  \item[\textbf{19.}] （1）证明见解析；（2）$60^\circ$。
  \item[\textbf{20.}] （1）$a_n=2^{7-n}$；（2）证明见解析。
  \item[\textbf{21.}] （1）$f(x)=-2x^3+3x^2$；（2）$0<m\leqslant\dfrac12$。
\end{itemize}
